% Solutions of 26-spanning-trees.
\section*{Chapter \ref{ch:b1:spanning-trees} --- Minimum Spanning Trees and Union--Find}

\begin{solution}{exo:b1:spanning-trees:1}
By increasing weight: $ac$ (1) kept, $bd$ (2) kept, $bc$ (3) kept, joining $\{a, c\}$ and $\{b, d\}$; $ab$ (4) and $cd$ (5)
close cycles; $de$ (6) kept. The tree $ac, bd, bc, de$ weighs 12.
\end{solution}

\begin{solution}{exo:b1:spanning-trees:2}
From $e$: $de$ (6) brings $d$; then $bd$ (2) brings $b$; then $bc$ (3) brings $c$; then $ac$ (1) brings $a$. The order is
$e, d, b, c, a$, and the tree is the same, of weight 12.
\end{solution}

\begin{solution}{exo:b1:spanning-trees:3}
\code{union(0, 1)}: 1 under 0, \omterm{def:b1:spanning-trees:rank}{rank} of 0 becomes 1. \code{union(2, 3)}: 3 under 2, \omterm{def:b1:spanning-trees:rank}{rank} of 2 becomes 1. \code{union(1,
3)}: the roots 0 and 2 both have \omterm{def:b1:spanning-trees:rank}{rank} 1, so 2 goes under 0, whose \omterm{def:b1:spanning-trees:rank}{rank} becomes 2. \code{union(4, 5)}: 5 under 4, \omterm{def:b1:spanning-trees:rank}{rank} 1.
\code{union(5, 0)}: the root 4 has \omterm{def:b1:spanning-trees:rank}{rank} 1, the root 0 \omterm{def:b1:spanning-trees:rank}{rank} 2, so 4 goes under 0. Parents: $0 \to 0$, $1 \to 0$, $2 \to
0$, $3 \to 2$, $4 \to 0$, $5 \to 4$; \omterm{def:b1:spanning-trees:rank}{ranks} 2 for 0, 1 for 2 and 4, 0 for the others.
\end{solution}

\begin{solution}{exo:b1:spanning-trees:4}
Yes. Each edge outside the tree is the heaviest of a cycle: $ab$ (4) of $a, b, c$; $cd$ (5) of $c, d, b$; $ce$ (7) of $c,
e, d, b$. By the cycle property none of them is in a \omterm{def:b1:spanning-trees:mst}{minimum spanning tree}, and the four edges left are exactly the
given tree.
\end{solution}

\begin{solution}{exo:b1:spanning-trees:5}
An edge whose two ends are both visited is rejected, even when they lie in different trees of the \omterm{def:b1:graph-traversals:free-tree}{forest}: two groups of
places that grew apart are never joined. The edges kept close no cycle, since one end is always new, so they form a
\omterm{def:b1:graph-traversals:free-tree}{forest}; 21 edges on 30 places leave $30 - 21 = 9$ trees. Testing ``same tree?'' needs \omterm{def:b1:spanning-trees:union-find}{union--find}, not a flag.
\end{solution}

\begin{solution}{exo:b1:spanning-trees:6}
On equal \omterm{def:b1:spanning-trees:rank}{ranks} it raises the \omterm{def:b1:spanning-trees:rank}{rank} of \code{ry}, the root that has just become a child, instead of the new root \code{rx}.
Roots keep \omterm{def:b1:spanning-trees:rank}{rank} 0, so every union is between equal \omterm{def:b1:spanning-trees:rank}{ranks} and links the root of \code{y} under the root of \code{x}. In
\code{union(i, i - 1)}, \code{i} is a fresh singleton: the whole chain goes under it, one level deeper each time, and
the finds of \code{i} and \code{i - 1}, at the top, compress nothing. After 999 unions the chain has height 999.
\end{solution}

\begin{solution}{exo:b1:spanning-trees:7}
\omcode{code/bachelor-1/26-spanning-trees/exercises/second_best.py}{12}{23}{A second-best \omterm{def:b1:spanning-trees:mst}{spanning tree}}
It is a known fact, which the test checks on small \omterm{def:b1:graph-traversals:graph}{graphs} with distinct weights, that some second-best tree differs from
the \omterm{def:b1:spanning-trees:mst}{minimum spanning tree} by a single exchange. The program tries the $(n - 1)(m - n + 1)$ exchanges and checks each
candidate with \omterm{def:b1:spanning-trees:union-find}{union--find}: $O(n m \cdot n\, \alpha(n))$, enough for small \omterm{def:b1:graph-traversals:graph}{graphs}. The test compares with the second
smallest weight among all \omterm{def:b1:spanning-trees:mst}{spanning trees}.
\end{solution}

\begin{solution}{exo:b1:spanning-trees:8}
\omcode{code/bachelor-1/26-spanning-trees/exercises/maze.py}{12}{22}{A maze by random Kruskal}
The cells and the walls between neighbours form a connected grid \omterm{def:b1:graph-traversals:graph}{graph}; a random order of the walls is a random order of
distinct weights, and Kruskal's algorithm opens the walls of a \omterm{def:b1:spanning-trees:mst}{spanning tree}. Connected: every cell is reachable; no
cycle: by exactly one way. The seed makes the maze the same at every run.
\end{solution}

\begin{solution}{exo:b1:spanning-trees:9}
By induction over the operations. At the start every root has \omterm{def:b1:spanning-trees:rank}{rank} 0 and 1 element. A root's \omterm{def:b1:spanning-trees:rank}{rank} goes from $r$ to $r +
1$ only when another root of \omterm{def:b1:spanning-trees:rank}{rank} $r$ is put under it; each had at least $2^r$ elements, so the merged tree has at least
$2^{r+1}$. \omterm{def:b1:compilation-pipeline:compiler}{Linking} under a root of larger \omterm{def:b1:spanning-trees:rank}{rank} only adds elements. \omterm{def:b1:spanning-trees:rank}{Path compression} changes parents inside a tree,
never which elements a root's tree holds, nor any \omterm{def:b1:spanning-trees:rank}{rank}.
\end{solution}

\begin{solution}{exo:b1:spanning-trees:10}
Suppose $T_1 \neq T_2$ are \omterm{def:b1:spanning-trees:mst}{minimum spanning trees}, and let $e$ be the lightest edge in one but not the other, say in
$T_1$. Adding $e$ to $T_2$ closes a cycle, which has an edge $f$ not in $T_1$, since $T_1$ has no cycle. Then $f$ is in
$T_2$ but not $T_1$, so $w(f) > w(e)$, the weights being distinct and $e$ the lightest such edge. $T_2 + e - f$ is a
\omterm{def:b1:spanning-trees:mst}{spanning tree} lighter than $T_2$: a contradiction.
\end{solution}

\begin{solution}{exo:b1:spanning-trees:11}
Let $T$ be a \omterm{def:b1:spanning-trees:mst}{minimum spanning tree} with heaviest edge $e$, and $B$ any \omterm{def:b1:spanning-trees:mst}{spanning tree}. Removing $e$ splits $T$ into two
parts, a \omterm{def:b1:spanning-trees:cut}{cut}; $B$ is connected, so it has an edge $f$ across that \omterm{def:b1:spanning-trees:cut}{cut}. If $w(f) < w(e)$, $T - e + f$ would be a lighter
\omterm{def:b1:spanning-trees:mst}{spanning tree}; so $w(f) \geq w(e)$, and the heaviest edge of $B$ weighs at least $w(e)$. For the converse, take a
triangle with $ab$ of weight 1 and $bc$, $ca$ of weight 2: the tree $bc, ca$ has heaviest edge 2, as small as possible,
but weighs 4 against 3.
\end{solution}

\begin{solution}{exo:b1:spanning-trees:12}
\omcode{code/bachelor-1/26-spanning-trees/exercises/boruvka.py}{11}{28}{Borůvka's method}
With weights made distinct by comparing the tuples, each picked edge is the only lightest edge leaving its tree, so it
is in the \omterm{def:b1:spanning-trees:mst}{minimum spanning tree} by the \omterm{def:b1:spanning-trees:cut}{cut} property, and the picked edges close no cycle. Each tree is merged with at
least one other in a round, so the number of trees at least halves: at most $\log_2 n$ rounds, each $O(m\, \alpha(n))$.
On Brenmoor it takes 3 rounds and finds 27\,210 m; the test compares it with Kruskal's algorithm on random \omterm{def:b1:graph-traversals:graph}{graphs}.
\end{solution}

\begin{solution}{pb:b1:spanning-trees:1}
\textbf{1.} Each edge that closes no cycle joins two components, so $n - 1$ of them leave $n - (n - 1) = 1$ component:
the \omterm{def:b1:graph-traversals:graph}{graph} is connected and acyclic, a \omterm{def:b1:graph-traversals:free-tree}{free tree} on all the \omterm{def:b1:graph-traversals:graph}{vertices}.

\textbf{2.} The 128 others contain a cycle; with 8 edges on 9 places, a set with a cycle leaves some place out.

\textbf{3.} $\binom{53}{29} \approx 7.8 \times 10^{14}$ sets, each checked with \omterm{def:b1:spanning-trees:union-find}{union--find}: years of computing.

\textbf{4.} \code{all} stops at the first \code{False}: a failed union means that the edge closes a cycle, and the set is
not a tree whatever the other edges.

\textbf{5.} That Kruskal's algorithm found an optimal tree on this piece; not that it does on every \omterm{def:b1:graph-traversals:graph}{graph}, which is the
proof's job.

\textbf{6.} A \omterm{def:b1:graph-traversals:connected}{connected graph} has at least $m - n + 1$ independent cycles: here $11 - 9 + 1 = 3$, so at least three
cycles.

\textbf{7.} 27\,210 m, along 29 roads.

\textbf{8.} Both compute \omterm{def:b1:spanning-trees:mst}{minimum spanning trees}, and all \omterm{def:b1:spanning-trees:mst}{minimum spanning trees} have the smallest weight, by
definition.

\textbf{9.} No: 53 roads have 47 distinct lengths. Two \omterm{def:b1:spanning-trees:mst}{minimum spanning trees} may then exist, of the same weight; the
program prints one.

\textbf{10.} 29 pops add a place; the other 15 take off edges whose two ends were already in the tree. Each road is
pushed at most twice, once from each end.

\textbf{11.} Either: Kruskal's with \omterm{def:b1:spanning-trees:union-find}{union--find} costs $O(m \log n)$ for the sort, and Prim's with a heap $O(m \log n)$;
the \omterm{def:b1:graph-traversals:graph}{graph} is sparse, so Prim's $O(n^2)$ \omterm{def:b1:pointers-and-arrays:array}{array} version would be the wrong one.

\textbf{12.} It is the lightest road across the \omterm{def:b1:spanning-trees:cut}{cut} that separates the South School from the rest: the school's other
roads are longer still, and the longer roads elsewhere close cycles.

\textbf{13.} Labels: each of the 99 unions scans the 100 elements, 9900 steps. \omterm{def:b1:pointers-and-arrays:pointer}{Pointers}: after $i$ unions the chain has
height $i$; the $i$-th union climbs $i - 1$ steps in \code{find(0)} and the following find $i$ steps:
$\sum_{i=1}^{99} (2i - 1) = 99^2 = 9801$.

\textbf{14.} Each fresh element goes under the root, and 0 hangs from the root after the first compression: each union
and each find climbs at most one \omterm{def:b1:pointers-and-arrays:pointer}{pointer}, $2(n - 2) + 1$ in all.

\textbf{15.} The root reaches \omterm{def:b1:spanning-trees:rank}{rank} 1 at the first union; every later union links a singleton of \omterm{def:b1:spanning-trees:rank}{rank} 0 under it.

\textbf{16.} Random unions mostly link a small tree under a large one, which does not raise the height; trees of equal
large \omterm{def:b1:spanning-trees:rank}{rank} rarely meet.

\textbf{17.} Merge the singletons in pairs, then the pairs in pairs, and so on: always two trees of equal \omterm{def:b1:spanning-trees:rank}{rank}, whose
height grows by one at each level, $\log_2 n$ levels in all.

\textbf{18.} Computing the new heights would cost more than the find; \omterm{def:b1:spanning-trees:rank}{ranks} stay upper bounds on heights, which is all
the proof of the logarithmic bound needs.

\textbf{19.} Let $e$ be the heaviest edge on the path from $a$ to $b$ in a \omterm{def:b1:spanning-trees:mst}{minimum spanning tree} $T$. Removing $e$ splits
$T$ into two parts, $a$ on one side and $b$ on the other. Any path from $a$ to $b$ \omterm{def:b1:spanning-trees:cut}{crosses} this \omterm{def:b1:spanning-trees:cut}{cut} at some edge $f$, and
$w(f) \geq w(e)$, or $T - e + f$ would be a lighter \omterm{def:b1:spanning-trees:mst}{spanning tree}. So every path has an edge at least as heavy as $e$.

\textbf{20.} Every route from the Town Square to the Hospital uses a road of at least 1240 m, and the route through the
tree uses none longer.

\textbf{21.} By Dijkstra's algorithm with $\max$ in place of $+$, which is correct for the same reason as the widest-path
variant of \cref{exo:b1:shortest-paths:8}; the two answers agree by question 19.

\textbf{22.} The longest tree road leads to the South School, a \omterm{def:b1:induction-and-termination:rooted-tree}{leaf} of the tree: removing it isolates one place. The next
one, Post Office--College, \omterm{def:b1:spanning-trees:cut}{cuts} off three places.

\textbf{23.} Not by itself: removing the longest roads peels off outlying places one by one rather than splitting the
town into balanced districts. It is the right method when clusters are separated by gaps, which a town's roads are not.

\textbf{24.} \emph{The minimum fibre length of Brenmoor} is 27\,210 m, computed by Kruskal's and Prim's algorithms alike.
By the bottleneck theorem (question 19 and \cref{exo:b1:spanning-trees:11}), its longest road, 1540 m, is the smallest
possible longest road of any network joining all the places, and the tree's paths are minimax routes.
\end{solution}
